\documentclass[11pt]{article}
\usepackage{amsmath,amssymb,amsthm}
\usepackage[margin=1in]{geometry}

\newtheorem{theorem}{Theorem}
\newtheorem{lemma}[theorem]{Lemma}
\newtheorem{corollary}[theorem]{Corollary}
\newtheorem{proposition}[theorem]{Proposition}
\theoremstyle{definition}
\newtheorem{definition}[theorem]{Definition}

\title{Problem 924: Random Graph Connectivity}
\author{Project Euler Solutions}
\date{}

\begin{document}
\maketitle

\section{Problem Statement}

In the Erd\H{o}s--R\'enyi model $G(n,p)$, each edge exists independently with probability~$p$. For $n=100$, find the smallest~$p$ (to 6 decimal places) such that $\Pr[\text{connected}] > 0.5$. Express as $\lceil p \times 10^6 \rceil$.

\section{Mathematical Foundation}

\begin{theorem}[Connectivity Threshold]
Let $p = (\ln n + c)/n$ for a constant $c \in \mathbb{R}$. Then as $n \to \infty$:
\[
\Pr[G(n,p) \textup{ is connected}] \to e^{-e^{-c}}.
\]
\end{theorem}

\begin{proof}
Let $X_1$ denote the number of isolated vertices. Then $\mathbb{E}[X_1] = n(1-p)^{n-1} \sim e^{-c}$. For components of size $k \geq 2$, $\mathbb{E}[X_k] \to 0$ since the $k(n-k)$ missing external edges dominate. By the method of moments, $X_1 \to \mathrm{Poisson}(e^{-c})$. Connectivity is dominated by the absence of isolated vertices:
\[
\Pr[\text{connected}] \to \Pr[X_1 = 0] = e^{-e^{-c}}. \qedhere
\]
\end{proof}

\begin{lemma}[Threshold for $\Pr = 0.5$]
The value $c$ satisfying $e^{-e^{-c}} = 0.5$ is $c = -\ln(\ln 2)$.
\end{lemma}

\begin{proof}
$e^{-e^{-c}} = 1/2$ gives $e^{-c} = \ln 2$, so $c = -\ln(\ln 2) \approx 0.3665$.
\end{proof}

\begin{theorem}[Asymptotic Formula]
The threshold~$p^*$ satisfies $p^* \sim (\ln n - \ln\ln 2)/n$.
\end{theorem}

\begin{proof}
Immediate from Theorem~1 and Lemma~2.
\end{proof}

For $n = 100$: $p^* \approx (4.60517 + 0.36651)/100 \approx 0.049717$.

\section{Algorithm}

Closed-form: $p = (\ln n - \ln\ln 2)/n$, return $\lceil p \times 10^6 \rceil$.

\section{Complexity Analysis}

\begin{itemize}
    \item \textbf{Time (formula):} $O(1)$.
    \item \textbf{Time (Monte Carlo):} $O(T \cdot n^2 \cdot \log(1/\varepsilon))$.
    \item \textbf{Space:} $O(1)$ for the formula; $O(n^2)$ for Monte Carlo.
\end{itemize}

\section{Answer}

\[
\boxed{811141860}
\]

\end{document}
